\chapter{討論與結論}
\label{ch:conclusion}

\section{系統優點與限制}
\label{sec:strengths-limitations}

\textbf{優點。}系統的設計特點已於\secref{sec:contributions}說明，\chapref{ch:experiments}的實驗進一步證實了自動化的上游步驟在未參與開發的比賽上依然有效。賽事影片分段與擊球時刻偵測在測試集上的 F1 分別為 98.66\% 與 93.14\%，與開發集的 99.34\% 與 90.70\% 相當。

\textbf{限制。}

\begin{enumerate}
  \item \textbf{依賴對轉播形式的假設。}多數規則建立在 BWF 轉播的慣例上，固定俯拍主機位、球場大致位於畫面中央偏下，且只處理單打。在符合假設的影片上準確率高，但假設一旦不成立，模組可能直接失效而非逐漸退化，例如手持或多機快切的影片無法分段，機位偏移可能使場地顏色分割整體失效。資料集 16 支影片中有 15 支來自 BWF，這類情況在實驗中幾乎沒有機會出現，此部分泛化能力尚未驗證。
  \item \textbf{規則有過擬合的跡象。}分段、場地偵測與擊球偵測都包含許多依開發集調整的門檻與規則。分段與擊球偵測在測試集上的表現與開發集相近，但場地偵測的命中率由開發集的 82.14\% 降至測試集的 72.22\%，失敗幾乎都來自開發時未遇過的底線廣告字樣。規則愈多，遇到新場館或新轉播風格時愈可能出現未預期的交互作用。本系統以規則取代標記訓練，避開了\chapref{ch:related-work}指出的標記成本，但依開發集調出的門檻同樣可能隨轉播形式失準。在資料集以 BWF 轉播為主的情況下，尚無法證明此路線比學習式方法更能泛化。
  \item \textbf{無法察覺自身失敗。}場地偵測在三場四角全錯的比賽中仍正常輸出，分段也可能整段漏掉回合而不留任何警示。系統缺乏信心度判斷，目前只能靠使用者回看時發現並手動修正。
  \item \textbf{外部依賴與處理速度。}比分辨識與賽評依賴 Gemini API，前者受速率上限影響，約占處理時間的 28\%。處理時間約需回合總長的 2.7 倍，且未含音訊、身分對應與賽評。若要用於延遲直播，仍需改為逐回合處理，並降低對 API 的依賴。
  \item \textbf{音訊支線的證據有限。}歡呼偵測只以四場開發比賽與一場外部驗證評估，不足以涵蓋不同轉播商、場館、觀眾規模與混音方式。該外部驗證場之後又參與了強度與聚合方法的開發，而這兩項的人工比較樣本少、只有單一標註者，因此只足以支持第一版的方法選擇。此外，歡呼強度是單場內的相對排名，不能跨比賽比較。歡呼也不等同回合的技術品質，可能同時受比分、球星或爭議影響，加上精彩度分數的 0.5／0.5 權重沒有人工排序真值支持，只能作為補充的排序訊號。
  \item \textbf{缺乏使用者評估。}介面與賽評目前只經開發者自行使用，沒有對教練、球員等目標使用者進行可用性測試，本報告也未量化評估生成賽評的品質，系統實際能減少多少人工逐拍標記的工作仍待驗證。
\end{enumerate}

\section{結論}
\label{sec:conclusion}

本專題建立了一套全自動的羽球單打轉播分析系統。系統輸出每個主機位比賽片段的起訖時間與比分、每一拍的擊球影格、擊球者與球種，以及雙方球員的骨架、球場範圍與羽球軌跡，並據此產生精彩片段排序與賽評。既有方法所依賴的比賽片段起訖、球場範圍與擊球時刻，以及既有系統未曾直接讀取的計分板比分，在本系統中都已自動取得。在未參與開發的比賽上，賽事影片分段與擊球時刻偵測的 F1 分別達 98.66\% 與 93.14\%，管線已能在真實轉播影片上從頭到尾自動完成。最弱的一環是場地偵測，測試集的命中率為 72.22\%，未命中的角幾乎都源於底線旁的廣告字樣。

綜合各模組的實驗結果，上游步驟自動化之後，球種資料的瓶頸落在兩處。一是 BST 本身，將擊球幀換成真值後準確率只提升 1.16 個百分點，未看過的 3 場中 202 拍錯誤有 146 拍來自 BST。二是擊球時刻偵測的召回率，漏掉的拍無從分類，使端到端準確率比單看分類低約 8 個百分點。此外，系統無法察覺自身的失敗，場地偵測在三場四角全錯的比賽中仍正常輸出，錯誤會靜默傳入下游，目前只能靠人工回看發現。因此，除了提升單一模組的準確率，讓系統能判斷自身輸出是否可信同樣重要。

未來工作將針對上述限制展開。首先，蒐集非 BWF 轉播與更多場館的影片，驗證並改善泛化能力，同時為各模組加入信心度估計，在輸出可疑時提示使用者回看。其次，降低對外部 API 的依賴並改為逐回合處理，實現延遲直播的應用。此外，為比分辨識與球員身分對應建立真值並量化評估，擴充歡呼偵測資料，並建立精彩度的人工排序真值。最後，對教練與球員進行使用者研究，量化評估系統與生成賽評的實際效益。
